\section{Conclusion}
\label{sec:conclusion}

This paper extends rationale-based distillation by introducing a boundary-aware rationale selection pipeline. Following the original training framework, rationales generated by a teacher LLM are used as auxiliary supervision for a compact student model. Unlike the original setting, however, our method does not directly use all generated rationales. It first groups candidate rationales by input example, estimates label confidence through weighted voting, assigns difficulty bands, and selects rationales using source compatibility, grounding, answer support, brevity, and agreement signals.

The main proposed variant, Boundary-Mix, keeps one high-quality rationale from easy and boundary examples while removing hard examples with unstable candidate agreement. Experimental results show that this selection strategy improves the original rationale-distilled baseline on both evaluated datasets. Boundary-Mix improves ESNLI from 83.93\% to 84.16\% and CommonsenseQA from 60.36\% to 62.41\%, while keeping the same compact student-model setting. The boundary-family comparison further shows that different selection rules have different effects: Boundary-Bridge gives the strongest ESNLI result, while Boundary-Mix gives the most consistent completed result across datasets.

Overall, the study shows that rationale-based distillation should consider not only how to generate rationales, but also how to select rationales that are suitable for learning. Generated explanations can be useful supervision, but they can also be incomplete, unstable, or weakly grounded. Boundary-aware rationale selection provides a transparent and reproducible way to reduce noisy rationale supervision before student training. Future work can improve this direction by learning source priors from development data, replacing token-overlap grounding with semantic verification, and evaluating the method on additional reasoning tasks and low-resource settings.
